%% ============================================================
%% BÖLÜM 10 — Stokastik Süreçlere Giriş  [OPSİYONEL]
%% Olasılık Teorisi ve Matematiksel İstatistik
%% ============================================================
\chapter{Stokastik Süreçlere Giriş}
\label{chp:stokastik}

\bolumalinti{%
  The theory of stochastic processes is really a theory of families of random variables
  indexed by a parameter set.
}{Joseph Doob}

\begin{kazanimlar}
Bu bölümün sonunda okuyucu:
\begin{itemize}
  \item Stokastik süreç, örneklem yolu, durumlar uzayı kavramlarını tanımlar.
  \item Markov zincirlerini tanımlayıcı özellikleriyle açıklar; geçiş matrislerini ve Chapman-Kolmogorov denklemini uygular.
  \item Süreç sınıflarını (geri dönüşlü, geçici, ergodik, dönemsel) belirler ve durağan dağılımı hesaplar.
  \item Poisson sürecini aksiyomlardan türetir; süreçler arası ilişkileri tanımlar.
  \item Brownian hareketi tanımlar ve temel özelliklerini (bağımsız artışlar, süreklilik, Gaussian dağılım) listeler.
\end{itemize}
\end{kazanimlar}

\begin{motivasyon}
Tek bir rasgele değişken statik bir belirsizliği modellerken, \emph{stokastik süreç} zamanla evrilen belirsizliği modelleyen rasgele değişken ailesidir. Hisse senedi fiyatları, kuyruklar, epidemi yayılımı, iklim değişkenliği — bunların tümü stokastik süreçtir.

\medskip
\noindent\textbf{Köprü:} Bölüm~\ref{chp:martingaller}'daki filtrasyon ve uyum kavramları doğrudan burada kullanılacak. Markov zincirlerindeki geçişler aslında $\sigma$-cebirlerinin monoton dizisi anlamında "bilgi artışı" kavramıdır.
\end{motivasyon}

%% ============================================================
\section{Temel Tanımlar}
\label{sec:stokastik_temel}
%% ============================================================

\begin{tanim}[Stokastik süreç]\label{def:stokastik_surec}
$(\Omega,\mathcal{F},\Prob)$ olasılık uzayı ve $T$ indeks kümesi olsun. \textbf{Stokastik süreç} \emph{(stochastic process)}, $\{X_t : t\in T\}$ rasgele değişken ailesidir. $T\subseteq\mathbb{Z}$ ise \textbf{ayrık zamanlı}, $T\subseteq[0,\infty)$ ise \textbf{sürekli zamanlı} denir. $S = X_t(\Omega)$ kümesine \textbf{durum uzayı} \emph{(state space)} denir.
\end{tanim}

\begin{tanim}[Örneklem yolu]\label{def:orneklem_yolu}
Sabit $\omega\in\Omega$ için $t\mapsto X_t(\omega)$ fonksiyonuna \textbf{örneklem yolu} \emph{(sample path)} veya \textbf{yörünge} \emph{(trajectory)} denir.
\end{tanim}

\begin{ornek}[Basit rasgele yürüyüş]\label{ex:basit_yuruyus}
$\xi_1,\xi_2,\ldots\iid$ ile $\Prob(\xi_i=+1)=p$, $\Prob(\xi_i=-1)=1-p$. $S_0=0$, $S_n=\sum_{k=1}^n\xi_k$. $\{S_n:n\geq0\}$ tam sayı değerli ayrık zamanlı stokastik süreçtir.
\end{ornek}

%% ============================================================
\section{Markov Zincirleri}
\label{sec:markov_zincirleri}
%% ============================================================

\begin{kesif}
Basit rasgele yürüyüşte $S_{n+1}$'in dağılımı, tüm geçmiş $S_0,S_1,\ldots,S_n$ yerine yalnızca $S_n$'e bağlıdır. Bu "geçmişten bağımsızlık" özelliği Markov zincirlerin özüdür.
\end{kesif}

\begin{tanim}[Markov zinciri]\label{def:markov_zinciri}
Ayrık durum uzaylı ($S$ sayılabilir) ayrık zamanlı süreç $\{X_n\}$, aşağıdaki \textbf{Markov özelliğini} \emph{(Markov property)} sağlıyorsa Markov zinciri adını alır:
\[
  \Prob(X_{n+1}=j \mid X_0=i_0,\ldots,X_{n-1}=i_{n-1},X_n=i)
  = \Prob(X_{n+1}=j\mid X_n=i).
\]
\end{tanim}

\begin{tanim}[Geçiş olasılıkları]\label{def:gecis_olasiliklari}
\textbf{Tek adım geçiş olasılığı}: $p_{ij} = \Prob(X_{n+1}=j\mid X_n=i)$.
\textbf{Geçiş matrisi}: $P = (p_{ij})_{i,j\in S}$; $p_{ij}\geq0$, $\sum_j p_{ij}=1$ (stokastik matris).
\end{tanim}

\begin{teorem}[Chapman-Kolmogorov]\label{thm:chapman_kolmogorov}
$p_{ij}^{(m+n)} = \sum_{k\in S} p_{ik}^{(m)}p_{kj}^{(n)}$, yani $P^{(m+n)} = P^{(m)}P^{(n)} = P^{m+n}$.
\end{teorem}

\begin{proof}
Koşullu olasılığı $X_m$ üzerinden toplayalım:
\[
  p_{ij}^{(m+n)} = \Prob(X_{m+n}=j\mid X_0=i)
  = \sum_k \Prob(X_{m+n}=j\mid X_m=k)\Prob(X_m=k\mid X_0=i)
  = \sum_k p_{kj}^{(n)}p_{ik}^{(m)}.
\]
\end{proof}

\subsection{Sınıf Yapısı}

\begin{tanim}[Erişilebilirlik ve sınıf]\label{def:erisim}
$j$ durumu $i$'den \textbf{erişilebilir} \emph{(accessible)}: $p_{ij}^{(n)}>0$ olan $n\geq0$ varsa; $i\leftrightarrow j$ ise ikisi \textbf{iletişimli} \emph{(communicating)}. İletişimli durumlar kümesi \textbf{sınıf} \emph{(class)} oluşturur. Her durum bir sınıfta. Tüm durumlar tek sınıf ise zincir \textbf{geri dönülemez} \emph{(irreducible)}.
\end{tanim}

\begin{tanim}[Geri dönüşlü ve geçici]\label{def:geri_donusluluk}
Durum $i$ için $f_i = \Prob(\exists n\geq1: X_n=i\mid X_0=i)$ geri dönüş olasılığı.
\begin{itemize}
  \item $f_i=1$: \textbf{geri dönüşlü} \emph{(recurrent)}.
  \item $f_i<1$: \textbf{geçici} \emph{(transient)}.
\end{itemize}
\end{tanim}

\begin{teorem}[Geri dönüşlülük kriteri]\label{thm:geri_donus_krit}
$\sum_{n=0}^\infty p_{ii}^{(n)} = \infty \iff i$ geri dönüşlü.
\end{teorem}

\begin{proof}
$R = \sum_{n=1}^\infty \mathbf{1}[X_n=i]$ geri dönüş sayısı. $\Beklenti_i[R]=\sum_{n=1}^\infty p_{ii}^{(n)}$. Geometrik dağılım argümanıyla $\Beklenti_i[R]=f_i/(1-f_i)$: $f_i=1$ iff $\Beklenti_i[R]=\infty$ iff seri ıraksar.
\end{proof}

\begin{tanim}[Periyot]\label{def:periyot}
$d(i) = \gcd\{n\geq1: p_{ii}^{(n)}>0\}$ durum $i$'nin \textbf{periyodu} \emph{(period)}. $d(i)=1$ ise \textbf{aperiyodik} \emph{(aperiodic)}.
\end{tanim}

\subsection{Durağan Dağılım}

\begin{tanim}[Durağan dağılım]\label{def:duraganlik}
Olasılık vektörü $\pi = (\pi_j)_{j\in S}$ ($\pi_j\geq0$, $\sum\pi_j=1$), $\pi P=\pi$ koşulunu sağlıyorsa \textbf{durağan dağılım} \emph{(stationary distribution)} denir.
\end{tanim}

\begin{teorem}[Ergodik teorem]\label{thm:ergodik}
$\{X_n\}$ geri dönülemez, geri dönüşlü, aperiyodik Markov zinciri olsun. O zaman:
\begin{enumerate}[(a)]
  \item Yegâne durağan dağılım $\pi$ vardır: $\pi_j = 1/m_j$ ($m_j$ = ortalama geri dönüş süresi).
  \item Her $i,j$: $p_{ij}^{(n)}\to\pi_j$ ($n\to\infty$).
  \item Her $f:S\to\mathbb{R}$ sınırlı için zaman ortalaması olasılık ortalamasına yakınsar:
  \[
    \frac{1}{n}\sum_{k=0}^{n-1}f(X_k) \asto \sum_j f(j)\pi_j.
  \]
\end{enumerate}
\end{teorem}

\begin{ispatiskele}
Geri dönüş teoremi, yenileme teorisine dayanır. $m_j=\Beklenti_j[T_j]$ (ilk geri dönüş zamanı). Ergodik limit, büyük sayılar kanunu analog. Tam ispat: Norris (1997) §1.8.
\end{ispatiskele}

\begin{ornek}[İki durumlu zincir]\label{ex:iki_durumlu}
$S=\{0,1\}$, $P=\begin{pmatrix}1-\alpha&\alpha\\\beta&1-\beta\end{pmatrix}$. Durağan dağılım: $\pi_0=\beta/(\alpha+\beta)$, $\pi_1=\alpha/(\alpha+\beta)$. $\alpha,\beta>0$ ise ergodik.
\end{ornek}

\begin{ornek}[Gambler's ruin — Markov zinciri]\label{ex:gambler_markov}
$S=\{0,1,\ldots,N\}$; $p_{i,i+1}=p$, $p_{i,i-1}=q=1-p$, sınır durumları absorbe edici. Bu Markov zincirinde $0$ ve $N$ geri dönüşlü, diğerleri geçici. $i$'den başlayıp $N$'ye ulaşma olasılığı $\rho_i$:
\[
  \rho_i = \frac{1-(q/p)^i}{1-(q/p)^N} \quad (p\neq q), \qquad \rho_i = i/N \quad (p=q).
\]
\end{ornek}

%% ============================================================
\section{Poisson Süreci}
\label{sec:poisson_sureci}
%% ============================================================

\begin{tanim}[Sayma süreci]\label{def:sayma_sureci}
$N(t)$ ($t\geq0$) tam sayı değerli süreç; $N(0)=0$ ve $t\mapsto N(t)$ n.k.\ sağdan sürekli, artan ise \textbf{sayma süreci} \emph{(counting process)}.
\end{tanim}

\begin{tanim}[Poisson süreci]\label{def:poisson_sureci}
Sayma süreci $\{N(t)\}$, aşağıdakileri sağlıyorsa $\lambda>0$ hızlı \textbf{Poisson süreci} \emph{(Poisson process)} adını alır:
\begin{enumerate}[(i)]
  \item $N(0)=0$.
  \item \textbf{Bağımsız artışlar}: $0\leq s<t$, $N(t)-N(s)$ önceki artışlardan bağımsız.
  \item \textbf{Durağan artışlar}: $N(t+s)-N(s) \overset{d}{=} N(t)$.
  \item $\Prob(N(h)=1)=\lambda h+o(h)$, $\Prob(N(h)\geq2)=o(h)$.
\end{enumerate}
\end{tanim}

\begin{teorem}[Poisson dağılımı]\label{thm:poisson_dagilim}
$\{N(t)\}$ Poisson süreci ise $N(t)\sim\Pois(\lambda t)$.
\end{teorem}

\begin{proof}
$p_n(t)=\Prob(N(t)=n)$ olsun. İleriye denklem: bağımsız artışlar ve aksiyom (iv) ile
\[
  p_n(t+h) = p_n(t)(1-\lambda h) + p_{n-1}(t)\lambda h + o(h).
\]
$h\to0$: $p_n'(t) = -\lambda p_n(t)+\lambda p_{n-1}(t)$, $p_0(0)=1$, $p_n(0)=0$ ($n\geq1$). Çözüm: $p_n(t)=e^{-\lambda t}(\lambda t)^n/n!$.
\end{proof}

\begin{teorem}[Varış araları]\label{thm:varis_aralari}
$T_k$ ($k$-inci varış arası) i.i.d.\ $\Exp(\lambda)$.
\end{teorem}

\begin{proof}
$\Prob(T_1>t) = \Prob(N(t)=0)=e^{-\lambda t}$; dolayısıyla $T_1\sim\Exp(\lambda)$. Markov özelliği (bağımsız artışlar) $T_2,T_3,\ldots$'yı da bağımsız $\Exp(\lambda)$ yapar.
\end{proof}

\begin{teorem}[İnceleme — süperpozisyon]\label{thm:poisson_ozellikler}
$\{N_1(t)\}$ ve $\{N_2(t)\}$ bağımsız, sırasıyla $\lambda_1$ ve $\lambda_2$ hızlı Poisson süreçleri.
\begin{enumerate}[(a)]
  \item \textbf{Süperpozisyon:} $N(t)=N_1(t)+N_2(t)$ ise $(\lambda_1+\lambda_2)$ hızlı Poisson.
  \item \textbf{İnceleme:} Her olayı $p$ olasılıkla tut. Tutulan $\{M(t)\}$ ise $p\lambda$ hızlı Poisson.
  \item \textbf{Koşullu:} $N(t)=n$ verildiğinde $n$ varış zamanı $[0,t]$ üzerinde i.i.d.\ $\Uniform[0,t]$.
\end{enumerate}
\end{teorem}

\begin{proof}[Proof (a)]
Karakteristik fonksiyonla: $N_1(t)+N_2(t)$ bağımsız Poisson toplamı $\Pois((\lambda_1+\lambda_2)t)$.
\end{proof}

%% ============================================================
\section{Brownian Hareket}
\label{sec:brownian}
%% ============================================================

\begin{tarihselnot}
Robert Brown (1827) su altındaki polenlerin rastgele hareketini gözlemledi. Einstein (1905) diffüzyon denklemini türetti. Wiener (1923) matematiksel temeli kurdu — bu nedenle \emph{Wiener süreci} de denir.
\end{tarihselnot}

\begin{tanim}[Standart Brownian hareket]\label{def:brownian}
$\{W(t): t\geq0\}$ süreci aşağıdakileri sağlıyorsa \textbf{standart Brownian hareket} \emph{(standard Brownian motion / Wiener process)} adını alır:
\begin{enumerate}[(i)]
  \item $W(0)=0$ n.k.
  \item \textbf{Bağımsız artışlar:} $0\leq s<t$, $W(t)-W(s)$ öncesinden bağımsız.
  \item \textbf{Gaussian artışlar:} $W(t)-W(s)\sim\Normal(0,t-s)$.
  \item \textbf{Sürekli yollar:} $t\mapsto W(t)$ n.k.\ sürekli.
\end{enumerate}
\end{tanim}

\begin{teorem}[Brownian hareketin temel özellikleri]\label{thm:brownian_ozellik}
$\{W(t)\}$ standart Brownian hareket. O zaman:
\begin{enumerate}[(a)]
  \item $W(t)\sim\Normal(0,t)$.
  \item $\Cov(W(s),W(t))=\min(s,t)$.
  \item Yollar n.k.\ hiçbir noktada türevlenemez.
  \item $W$ bir martingaldir.
\end{enumerate}
\end{teorem}

\begin{proof}
(a) $W(t)=W(t)-W(0)\sim\Normal(0,t)$ tanımdan.

(b) $s\leq t$: $\Cov(W(s),W(t))=\Cov(W(s),W(s)+(W(t)-W(s)))=\Var(W(s))=s$ (bağımsız artışlar).

(c) Yol türevsizliği: Paley-Wiener-Zygmund teoremi (ileri konu, referans: Mörters-Peres).

(d) $\Beklenti[W(t)\mid\mathcal{F}_s]=W(s)$: $W(t)-W(s)\perp\mathcal{F}_s$, $\Beklenti[W(t)-W(s)]=0$.
\end{proof}

\begin{ornek}[Geometrik Brownian hareket]\label{ex:geometrik_brownian}
$S(t) = S_0\exp\!\bigl[(\mu-\tfrac{1}{2}\sigma^2)t + \sigma W(t)\bigr]$ formülü Black-Scholes modelinde hisse senedi fiyatını verir. $\log S(t)\sim\Normal(\log S_0+(\mu-\sigma^2/2)t,\,\sigma^2 t)$.
\end{ornek}

\begin{disiplinlerarasibaglan}
\textbf{Fizik:} Brownian hareket ısı difüzyonunu modelleyen $\partial u/\partial t = \frac{1}{2}\partial^2 u/\partial x^2$ denkleminin olasılıksal çözümüdür.
\textbf{Finans:} Black-Scholes opsiyon fiyatlaması geometrik Brownian harekete dayanır.
\textbf{Biyoloji:} Protein difüzyonu, nöral gürültü modelleri.
\end{disiplinlerarasibaglan}

%% ============================================================

\begin{mikroozet}
\begin{itemize}
  \item Markov zinciri: gelecek yalnızca şimdiye bağlı; Chapman-Kolmogorov $P^{m+n}=P^m P^n$.
  \item Durağan dağılım: $\pi P=\pi$; ergodik zincirde yegâne ve $p_{ij}^{(n)}\to\pi_j$.
  \item Poisson süreci: bağımsız+durağan artışlar; $N(t)\sim\Pois(\lambda t)$; varış araları $\Exp(\lambda)$.
  \item Brownian hareket: Gaussian bağımsız artışlar, sürekli türevsiz yollar, martingal.
\end{itemize}
\end{mikroozet}

\begin{ozyansima}
Markov özelliği "geçmişi unutan" sistemleri ne kadar iyi modelliyor? Hangisi daha zengin: Markov zincirleri mi yoksa martingaller mi? Bu ikisinin kesişimi nedir?
\end{ozyansima}

\begin{kopru}
\textbf{Nereden geldik:} Bölüm~\ref{chp:martingaller}'daki Doob yakınsama teoremi, geri dönüşlü Markov zincirlerin ergodik teoreminin özel durumudur.

\textbf{Nereye gidiyoruz:} Part II'de (Böl.~\ref{chp:guven_araliklari} vd.) bu süreçlerin parametrelerini veriden tahmin etmeye odaklanacağız. Markov zinciri Monte Carlo (MCMC) Bayesçi istatistikte doğrudan kullanılır.
\end{kopru}

%% ============================================================
%% ALISTIRMALAR
%% ============================================================

\cozumlualistirmalar

\begin{alistirma}[Al.10.1]{A}{Durağan dağılım}
Üç durumlu Markov zinciri $P=\begin{pmatrix}0.5&0.4&0.1\\0.3&0.5&0.2\\0.2&0.3&0.5\end{pmatrix}$. Durağan dağılımı bul.
\end{alistirma}
\alcozum{%
$\pi P=\pi$, $\sum\pi_i=1$. Denklem sistemi: $0.5\pi_0+0.3\pi_1+0.2\pi_2=\pi_0$ vb. Çözüm: $\pi=(0.303,0.404,0.293)$ (yaklaşık; tam çözüm için lineer denklem sistemi).%
}

\begin{alistirma}[Al.10.2]{B}{Poisson süperpozisyon}
Bir hastane acil birimine birbirinden bağımsız iki kaynaktan hasta geliyor: 112 araçlarından $\lambda_1=3$/saat, bireysel başvurulardan $\lambda_2=2$/saat. Bir saatte 10 veya daha fazla hasta gelme olasılığı nedir?
\end{alistirma}
\alcozum{%
Toplam süreç $\Pois(5)$. $\Prob(N(1)\geq10)=1-\Prob(N(1)\leq9)=1-\sum_{k=0}^9 e^{-5}5^k/k!\approx1-0.9682=0.0318$.%
}

\alistirmalar

\begin{alistirma}[Al.10.3]{A}{Brownian varyans}
$W(t)$ standart Brownian hareket. $\Beklenti[W(3)^2]$, $\Cov(W(2),W(5))$ ve $\Cor(W(2),W(5))$'i hesapla.
\end{alistirma}
\alcozum{%
$\Beklenti[W(3)^2]=\Var(W(3))=3$. $\Cov(W(2),W(5))=\min(2,5)=2$. $\Cor=2/\sqrt{2\cdot5}=2/\sqrt{10}$.%
}

\begin{alistirma}[Al.10.4]{B}{İki durumlu ergodiklik}
İki durumlu Markov zinciri $P=\begin{pmatrix}1-\alpha&\alpha\\\beta&1-\beta\end{pmatrix}$ ($\alpha,\beta\in(0,1)$). (a) Durağan dağılımı bul. (b) $p_{00}^{(n)}$'in kapalı formülünü türet. (c) $n\to\infty$ limitinin durağan dağılıma yakınsadığını doğrula.
\end{alistirma}
\alcozum{%
(a) $\pi_0=\beta/(\alpha+\beta)$, $\pi_1=\alpha/(\alpha+\beta)$. (b) Özdeğerler $1$ ve $1-\alpha-\beta$: $p_{00}^{(n)}=\pi_0+\pi_1(1-\alpha-\beta)^n$. (c) $|1-\alpha-\beta|<1$ olduğundan ikinci terim $\to0$; $p_{00}^{(n)}\to\pi_0$.%
}

\begin{alistirma}[Al.10.5]{C}{Gambler's ruin — genel}
$p\neq q$ için Gambler's Ruin'de $i$'den başlayıp $N$'ye ulaşma olasılığı $\rho_i=\frac{1-(q/p)^i}{1-(q/p)^N}$ olduğunu fark denklemiyle ispat et. Ortalama oyun uzunluğunu $i=N/2$, $p=q=1/2$ özel durumunda hesapla.
\end{alistirma}
\alcozum{%
Fark denklemi: $\rho_i=p\rho_{i+1}+q\rho_{i-1}$, $\rho_0=0$, $\rho_N=1$. Karakteristik denklem $pr^2-r+q=0$; kökler $1$ ve $q/p$. Genel çözüm $\rho_i=A+B(q/p)^i$; sınır koşullarıyla $\rho_i$ bulunur. $p=q=1/2$, $i=N/2$: ortalama süre $= i(N-i) = N^2/4$.%
}
